\section{Introduction}
\label{p1:intro}

A long constant run in independent fair signs starts with probability
\(2^{-L}\). Thus \(N\) possible starts have a bounded expected count
when \(L=\log_2N+O(1)\). A completely multiplicative sign sequence
has the same rare-event scale, but its distant windows share prime
signs. The resulting dependencies leave an explicit signature in the
Poisson approximation.

The first signature comes from pairs of windows in the ratio \(2:3\).
Their total contribution has size \(LN^{-2/3}\). We determine its
coefficient and then ask how accurately the whole counting law can be
recovered after the exact pair correction is retained. The answer has
optimal power \(4/3\): two pair interactions survive at fourth order,
whereas the total triple activity decays with power \(41/30\).
The gap \(41/30-4/3=1/30\) is enough, provided that the expansion
also controls configurations whose cardinality grows with \(N\).

\subsection{Model and principal results}
\label{p1:intro:model}

Choose independent fair signs \(f(\ell)\), indexed by primes, and set
\[
 f(n)=\prod_{\ell}f(\ell)^{v_\ell(n)}\qquad(n\ge1).
\]
The function \(f\) is completely multiplicative; in particular,
\(f(n^2)=1\). Fix \(C>0\). Throughout, \(N\ge2\) and \(L\ge1\)
are integers satisfying
\begin{equation}
 I_N=[N,2N)\cap\Z,\qquad
 |L-\log_2N|\le C,\qquad B=L+1,\qquad
 p=2^{-L},\qquad \lambda_N=Np.
 \label{p1:eq:critical}
\end{equation}
All logarithms without a subscript are natural. The nominal mean
\(\lambda_N\) remains in a compact subinterval of \((0,\infty)\),
but need not converge. Define
\begin{equation}
 J_x=\ind_{\{f(x-1)=-f(x),\
                 f(x)=f(x+1)=\cdots=f(x+L-1)\}},
 \qquad Z_N=\sum_{x\in I_N}J_x.
 \label{p1:eq:model}
\end{equation}
Both signs of the run are included. The left boundary condition counts
a long run once, at its start. A retained run may continue beyond
\(2N\), and a run starting before \(N\) is not counted.

Write
\[
 p_x=\E J_x,\qquad \mu_N=\E Z_N,\qquad
 \alpha_N=\Var(Z_N)-\mu_N.
\]
The exact marginals \(p_x\) need not equal \(p\).
For a signed measure \(\nu\) on \(\N=\{0,1,\ldots\}\), let
\(\|\nu\|_1=\sum_{k\ge0}|\nu\{k\}|\); total variation between
probabilities is \(\TV(P,Q)=\frac12\|P-Q\|_1\).
Set \(\Pi_\mu=\Pois(\mu)\) and \(\Delta=\delta_1-\delta_0\).
Products of measures are convolutions. For $r\ge1$, write
$(z)_r=z(z-1)\cdots(z-r+1)$, with $(z)_0=1$, and set
$\Delta^{*0}=\delta_0$. The quadratic reference is
\begin{equation}
 Q_{\mu,\alpha}
   =\Pi_\mu+\frac{\alpha}{2}\Pi_\mu*\Delta^{*2}.
 \label{p1:eq:Q}
\end{equation}
It has total mass one and, in the range of the theorem below, is
also nonnegative.

\begin{theorem}[Optimal quadratic approximation with exact parameters]
\label{p1:thm:main}
Fix \(C>0\). For every \(\varepsilon>0\), there are positive
constants \(c_{C,\varepsilon}\), \(C_{C,\varepsilon}\) and a threshold
\(N_0(C,\varepsilon)\) such that, for all \(N\ge N_0\) and every
\(L\) satisfying \eqref{p1:eq:critical},
\begin{equation}
 c_{C,\varepsilon}N^{-4/3-\varepsilon}
 \le \left\|\Law(Z_N)-Q_{\mu_N,\alpha_N}\right\|_1
 \le C_{C,\varepsilon}N^{-4/3+\varepsilon}.
 \label{p1:eq:main}
\end{equation}
Moreover,
\(\mu_N-\lambda_N=O_{C,\varepsilon}(N^{-1+\varepsilon})\),
and \(Q_{\mu_N,\alpha_N}\) is a probability measure for all
sufficiently large \(N\), uniformly in the same band.
\end{theorem}

The theorem determines the optimal power of the remainder, with an
arbitrary fixed power margin. It does not determine its logarithmic
factor or leading constant. The exact mean matters: replacing it by
\(\lambda_N\) can have a larger cost of order \(N^{-1+o(1)}\).
Proposition~\ref{p13:optimal:nominal} makes this distinction precise.

The pair calculation gives a full leading equivalent. If
\(r\in\{0,1,2\}\) is the residue of \(L\) modulo \(3\), then
\begin{equation}
 \alpha_N\sim
 \frac{r+2}{18\,2^{r/3}}\lambda_N^{5/3}L N^{-2/3}.
 \label{p1:eq:alpha-leading}
\end{equation}
Here and below, an equivalent with moving parameters means that the
ratio tends to one uniformly over the fixed critical band.

\begin{corollary}[Exact Poisson total variation scale]
\label{p1:cor:tv}
For every \(\varepsilon>0\), uniformly under \eqref{p1:eq:critical},
\begin{equation}
 \TV(\Law(Z_N),\Pi_{\lambda_N})
 =\frac{\alpha_N}{4}\|\Pi_{\lambda_N}*\Delta^{*2}\|_1
       +O_{C,\varepsilon}(N^{-1+\varepsilon}).
 \label{p1:eq:tv-main}
\end{equation}
Consequently,
\begin{equation}
 \TV(\Law(Z_N),\Pi_{\lambda_N})
 \sim\frac{r+2}{72\,2^{r/3}}\lambda_N^{5/3}
       \|\Pi_{\lambda_N}*\Delta^{*2}\|_1\,L N^{-2/3}.
 \label{p1:eq:tv-equiv}
\end{equation}
With \(\Pi_{\mu_N}\) as the reference, the first formula holds with
\(\mu_N\) in the difference norm and error
\(O_{C,\varepsilon}(N^{-4/3+\varepsilon})\); the equivalent is
unchanged.
\end{corollary}

Thus the power \(2/3\), the factor \(L\), and the leading coefficient
are determined for both specified Poisson means. No minimization over
all possible means is asserted.

\begin{example}[A fixed phase]
\label{p1:ex:unit}
Take \(N=2^{3k}\) and \(L=3k\). Then \(\lambda_N=1\), \(r=0\),
and \(\|\Pi_1*\Delta^{*2}\|_1=3/e\). Hence
\begin{equation}
 \TV(\Law(Z_N),\Pi_1)
       \sim\frac{1}{12e}\,L N^{-2/3}.
 \label{p1:eq:unit-example}
\end{equation}
The constant combines the arithmetic coefficient \(1/9\), the
second-difference norm \(3/e\), and the factor \(1/4\) from the
cumulant and total variation conventions.
\end{example}

A rational channel $a:b$ has coprime integers $a<b<2a$ and height $b$.
An explicit reference can retain the complete contributions of all
channels up to a fixed integer $h\ge3$. The following table distinguishes
the resulting approximation statements. All errors are measured in
total variation, uniformly in the fixed critical band; $h$ and every
margin $\varepsilon>0$ are fixed before $N\to\infty$.

\begin{center}
\small
\begin{tabularx}{\linewidth}
  {@{}p{.20\linewidth}p{.25\linewidth}X@{}}
\toprule
Reference & Parameters retained & Precision established here\\
\midrule
Poisson
& Nominal mean $\lambda_N$ or exact mean $\mu_N$
& Explicit equivalent of order $LN^{-2/3}$
  (Corollary~\ref{p1:cor:tv}).\\[3pt]
Exact quadratic
& $\mu_N,\alpha_N$
& Optimal power $4/3$, with arbitrary fixed power margins
  (Theorem~\ref{p1:thm:main}).\\[3pt]
Explicit quadratic
& Nominal mean and complete channel sums of heights $\le h$
& Explicit equivalent of order $LN^{-h/(h+1)}$
  (Theorem~\ref{p13:channels:hierarchy}).\\[3pt]
Exact Hermite
& $\mu_N,\alpha_N$
& Upper bound $O_{C,\varepsilon}(N^{-4/3+\varepsilon})$;
  optimality is not established
  (Corollary~\ref{p1:cor:compound}).\\
\bottomrule
\end{tabularx}
\end{center}

For $h=4$, the explicit quadratic reference retains the $2:3$ and $3:4$
channels and has optimal power $4/5$.
Theorem~\ref{p13:third:main} also determines the optimal power $41/30$
of the third factorial cumulant; its positive witness consists of
triples in the ratio $10:12:15$.

\subsection{How the first channel produces its scale}
\label{p1:intro:mechanism}

The interval selects the first interior ratio. If \(x\in I_N\),
then \(2x\notin I_N\), so the ratio \(1:2\) has no macroscopic
family of interior placements. The next ratio, \(2:3\), does.
For \(L=3m\), the windows at \(2t\) and \(3t\) contain the matched
values \(2(t+j)\) and \(3(t+j)\), \(0\le j<m\). For \(j\ne k\),
\[
 f(2(t+j))f(2(t+k))=f(3(t+j))f(3(t+k)).
\]
The canonical matching has rank \(m-1\) when there are no further
square relations. Its prescriptions are compatible because the
predecessor signs do not occur in these matches. Such a pair has
\[
 \P(J_{2t}=J_{3t}=1)=2^{-2L+m-1},
\]
rather than the independent weight \(2^{-2L}\).

Allowing the offset \(2y-3x\) to vary over a range of length
\(O(L)\) gives \(O(NL)\) placements. Their typical joint weight is
\(N^{-5/3}\), which explains the product
\[
             (NL)N^{-5/3}=LN^{-2/3}.
\]
This calculation identifies the mechanism. To obtain the coefficient,
one must sum the exact endpoint and parity conditions and bound pairs
with extra relations. Theorem~\ref{p1:arith:pairs} performs that
calculation. The complete finite \(2:3\) contribution matters:
its next term, of order \(N^{-2/3}\), is larger than the leading
\(3:4\) contribution of order \(LN^{-3/4}\). Successive channels
are therefore isolated by subtracting complete finite expressions.

\subsection{From arithmetic interactions to the whole law}
\label{p1:intro:strategy}
\label{p1:intro:dependencies}

The reference is first a sum of independent Bernoulli variables with
the exact probabilities \(p_x\). This retains every one-window effect,
including repeated indices. Each event is an affine condition on the
prime signs over \(\F_2\). Effective square-class spaces record its
constraints after local relations have been removed. A homogeneous
joint probability then dominates the exact rank weights, actual joint
probabilities and absolute centered activities, so a single geometric
estimate can control several terms in the expansion.

The finite stopping identity in Section~\ref{p1:stop:section} retains
a hereditary class of configurations with a computable probability.
It stops at the first excluded prefix and leaves the remaining
indicators inside a positive conditional measure. The resulting error
is controlled by the mass of that boundary. Fixed configurations are
handled by conditional arithmetic counts. A separate certificate
argument sums the large-cardinality tail. Bounded moments at every
fixed order do not replace this uniform estimate.

For
\[
 A_3=\sum_{\substack{T\subseteq I_N\\|T|=3}}
       \left|\E\prod_{x\in T}(J_x-p_x)\right|,
\]
Theorem~\ref{p13:transfer:reconstruction} gives
\[
 \|\Law(Z_N)-Q_{\mu_N,\alpha_N}\|_1
       \ll_{C,\varepsilon}N^\varepsilon A_3
                            +N^{-4/3+\varepsilon}.
\]
The arithmetic problem is therefore to place \(A_3\) below the
power \(4/3\). For a pair with sufficiently large shared rank, short
relations produce a common curve and force many integer points in a
logarithmic offset grid. A degree-uniform point estimate rules out this
rank for nonaligned pairs. Combining this restriction
with the explicit geometry of aligned triples gives
\(A_3\ll N^{-41/30+\varepsilon}\).
Taking a sufficiently small fixed margin yields
\[
 N^\varepsilon A_3
   \ll N^{-4/3-(1/30-2\varepsilon)},
 \qquad 0<\varepsilon<1/60.
\]
This is the arithmetic gain that closes the upper bound.

For the lower bound, two pair interactions have size
\(\alpha_N^2=N^{-4/3+o(1)}\). A fourth Charlier polynomial detects
the corresponding term in the true law. The proof keeps a
nonnegative compatible contribution and bounds the incompatible
part before applying independently proved fixed-moment estimates.
This establishes noncancellation for the specified quadratic
reference. The exact Hermite reference already contains a fourth-order
pair-pair term, so the same lower bound does not apply to it.

\subsection{Context and organization}

For independent signs, long runs and rare patterns have a classical
Poisson theory; see Arratia, Gordon, and Waterman \cite{agw1990}
and Godbole \cite{godbole1991}. In the multiplicative model,
Najnudel \cite[Lemma~3.1 and Proposition~6.3]{najnudel2020}
describes consecutive values through valuation relations and proves
almost-sure frequencies of fixed words. Fixed-length frequencies do
not supply the uniform rare-event estimates needed when
\(L\asymp\log N\). The binary completely multiplicative model is
also distinct from the squarefree-supported convention for a
Rademacher multiplicative function; the former is called extended
Rademacher in de la Bret\`eche, Wang, and Xu \cite{bretwngxu2026}.

Basak, Nath, and Zaharescu
\cite[Theorem~1]{basaknathzaharescu2024} prove related Poisson behavior
for chains of quadratic non-residues. Their statistic counts successful
disjoint blocks in moving windows of block indices, for almost all primes
in a short interval. The block length tends to infinity at most as fast
as four iterated logarithms of the prime scale. This differs from the
random run starts and critical length regime considered here.

Pouly \cite[Theorem~1.1]{pouly2026c} establishes a Poisson
approximation for the sign-resolved run field and already uses the \(2:3\)
interaction in a lower bound \cite[Section~3.6, (3.22)]{pouly2026c}.
Here the observable is the scalar count with both signs included.
Its finer approximation requires the complete pair coefficient,
control of the law below that coefficient, and a lower bound
for the residual after correction.

Poisson approximation for dependent indicators has a general
Chen--Stein framework \cite{agg1989}; asymptotic refinements include
\cite{barbour1987,barboureagleson1987}.
The reference \(Q_{\mu,\alpha}\) is explicit in Barbour and Eagleson
\cite[Theorem~1, p.~588]{barboureagleson1987}.
The Poisson--Charlier organization by factorial cumulants is classical
\cite[Sections~5 and~8.1]{bkn2014}.
\v{C}ekanavi\v{c}ius and Vellaisamy
\cite[Theorems~3.1--3.2]{cekanaviciusvellaisamy2015} distinguish
quadratic and exponential corrections in a locally dependent setting.
Their later lower bound
\cite[Theorem~2.1, (2.3)]{cekanaviciusvellaisamy2020}
exhibits the squared second-factorial-cumulant obstruction.
These local-dependence hypotheses do not cover the separated arithmetic
windows here. Our contribution is the arithmetic asymptotics of the
parameters and the uniform control and noncancellation needed for the
optimal remainder power.

The distance-equivalent theorem of
Chhaibi, Delbaen, M\'eliot, and Nikeghbali
\cite[Theorem~3.11]{chhaibietal2020} uses a diverging parameter and
regular Fourier residues; our Poisson means stay in a compact
subset of \((0,\infty)\). This restriction belongs to that theorem:
the finite bounds of \cite[Theorem~5.1]{bkn2014} also treat bounded
Poisson means once the required residue estimates are available.

The finite stopping identity is a specialization of the abstract
lace expansion \cite{zeilberger1997}. Its application here depends
on a retained class adapted to rational pairs and on arithmetic
bounds for its boundary. Those bounds, the explicit channel
coefficients and the noncancellation argument give the stated
precision for multiplicative runs.

This article and its \emph{Technical companion} form one work.
References beginning with T belong to the companion. The article
develops the exact finite expansion, states the arithmetic estimates
it consumes, and carries out reconstruction and optimality. The
companion proves the arithmetic and uniform counting estimates.
Every error margin, fixed order and channel height is chosen before
\(N\to\infty\); any growing range has a separate uniform estimate.
\label{p1:intro:notation}
